% block: matrix-function-definition-heading page:203 type:section
\subsection{{Definition and Properties of Matrix Functions}}\label{sec:matrix-function-properties}

% block: matrix-spectrum-introduction page:203 type:ordinary-paragraph
{The set of all the eigenvalues of matrix $A$ is called the spectrum of $A$. The definition of a function on the spectrum of a matrix is given as follows.}

% block: function-on-spectrum-definition page:203 type:definition
\Needspace{8\baselineskip}
\begin{Definition}{{function on the spectrum of a matrix}}\label{def:function-on-matrix-spectrum}
{Let $\lambda_1,\lambda_2,\ldots,\lambda_s$ be distinct eigenvalues of matrix $A$, and let $n_i$ be the size of a Jordan block with respect to eigenvalue $\lambda_i$ of $A$. If there exist the following values:}
\[f(\lambda_i),\frac{\mathrm{d}}{\mathrm{d}x}f(\lambda_i),\ldots,\frac{\mathrm{d}^{n_i-1}}{\mathrm{d}x^{n_i-1}}f(\lambda_i)\quad(i=1,2,\ldots,s),\]
{then $f(x)$ is called a function on the spectrum of matrix $A$.}
\end{Definition}

% block: matrix-function-derivative-notation page:203 type:ordinary-paragraph
{In what follows, we use the symbol $f^{(i)}(x)$ instead of $\frac{\mathrm{d}^i}{\mathrm{d}x^i}f(x)$.}

% block: inverse-function-spectrum-example page:203 type:example
\Needspace{7\baselineskip}
\begin{Example}\label{ex:inverse-on-jordan-spectrum}
{Let $A$ be matrix $J$ in Example~\ref{ex:jordan-block-diagonal-six}. Then the number of distinct eigenvalues is $s=2$ since $\lambda_1=\lambda_2=1$ and $\lambda_3=2$. The maximum sizes of Jordan blocks corresponding to $\lambda_1=\lambda_2=1$ and $\lambda_3=2$ are $n_1=2$ and $n_2=3$. The spectrum of matrix $A$ is $\{\lambda_1,\lambda_3\}$. We now give an example of a function on the spectrum of matrix $A$. Let $f(x)=x^{-1}$. Then $f(x)$ is a function on the spectrum of matrix $A$, because $f^{(1)}(x)=-x^{-2}$, $f^{(2)}(x)=2x^{-3}$ and there exist the values of the functions on the spectrum as follows: $f(\lambda_1)=1$, $f^{(1)}(\lambda_1)=-1$, $f(\lambda_3)=1/2$, $f^{(1)}(\lambda_3)=-1/4$, and $f^{(2)}(\lambda_3)=1/4$.}
\end{Example}

% block: jordan-matrix-function-transition page:203 type:ordinary-paragraph
{A matrix function of $A$ can be defined by using Jordan canonical form in Theorem~\ref{thm:jordan-canonical-form} and a function on the spectrum of a matrix in Definition~\ref{def:function-on-matrix-spectrum} as follows:}

% block: jordan-matrix-function-definition page:204 type:definition
\Needspace{9\baselineskip}
\begin{Definition}{{Matrix function}}\label{def:jordan-matrix-function}
{Let $f(x)$ be a function on the spectrum of matrix $A\in\mathbb{C}^{n\times n}$ and $A=ZJZ^{-1}$ be Jordan canonical form of $A$. Then, matrix function $f(A)=f(ZJZ^{-1})$ is defined as follows:}
\[f(A)=Z\operatorname{diag}(f(J_1),f(J_2),\ldots,f(J_p))Z^{-1}.\]
{Here $f(J_k)$ is the following $m_k\times m_k$ matrix:}
\[
f(J_k)=\begin{bmatrix}f(\lambda_k)&\frac{f^{(1)}(\lambda_k)}{1!}&\cdots&\frac{f^{(m_k-1)}(\lambda_k)}{(m_k-1)!}\\&f(\lambda_k)&\ddots&\vdots\\&&\ddots&\frac{f^{(1)}(\lambda_k)}{1!}\\&&&f(\lambda_k)\end{bmatrix}\in\mathbb{C}^{m_k\times m_k},
\]
{and $m_k\times m_k$ is the size of Jordan block corresponding to eigenvalue $\lambda_k$.}
\end{Definition}

% block: matrix-function-eigenvalues page:204 type:ordinary-paragraph
{It follows from Definition~\ref{def:jordan-matrix-function} that when $\lambda_k$ is an eigenvalue of matrix $A$, the eigenvalue of matrix function $f(A)$ is $f(\lambda_k)$.}

% block: matrix-function-diagonalizable-case page:204 type:ordinary-paragraph
{In particular, when matrix $A$ is diagonalizable, the Jordan blocks are diagonal matrices $D_k$ for $k=1,2,\ldots,p$. Then it follows from $f(A)=Z\operatorname{diag}(D_1,D_2,\ldots,D_p)Z^{-1}$ that $f(A)$ and $A$ have the same eigenvectors.}

% block: inverse-matrix-function-example page:204 type:example
\Needspace{9\baselineskip}
\begin{Example}\label{ex:inverse-jordan-matrix-function}
{Let $A$ be matrix $J$ in Example~\ref{ex:jordan-block-diagonal-six} and let $f(x)=x^{-1}$. Jordan canonical form of $A$ is given by $A=ZJZ^{-1}$, where $Z$ is the identity matrix. Then $f(A)=f(ZJZ^{-1})=f(J)$ and thus from $\lambda_1=\lambda_2=1$, $\lambda_3=2$ and Definition~\ref{def:jordan-matrix-function}, we have}
\[
\begin{aligned}
f(A)&=\left[\begin{array}{cc|c|ccc}f(\lambda_1)&\frac{f^{(1)}(\lambda_1)}{1!}&0&0&0&0\\0&f(\lambda_1)&0&0&0&0\\\hline0&0&f(\lambda_2)&0&0&0\\\hline0&0&0&f(\lambda_3)&\frac{f^{(1)}(\lambda_3)}{1!}&\frac{f^{(2)}(\lambda_3)}{2!}\\0&0&0&0&f(\lambda_3)&\frac{f^{(1)}(\lambda_3)}{1!}\\0&0&0&0&0&f(\lambda_3)\end{array}\right]\\
&=\left[\begin{array}{cc|c|ccc}1&-1&0&0&0&0\\0&1&0&0&0&0\\\hline0&0&1&0&0&0\\\hline0&0&0&2^{-1}&-4^{-1}&8^{-1}\\0&0&0&0&2^{-1}&-4^{-1}\\0&0&0&0&0&2^{-1}\end{array}\right].
\end{aligned}
\]
{It is easy to see that $f(A)$ is the inverse of matrix $A$, i.e., $A^{-1}$ was derived from the definition of the matrix function $f(A)$.}
\end{Example}

% block: square-root-exponential-jordan-example page:205 type:example
\Needspace{8\baselineskip}
\begin{Example}\label{ex:square-root-exponential-jordan}
{We consider other examples of matrix functions of $A$. Let}
\[A=\begin{bmatrix}-1&6&-9\\4&-4&12\\3&-5&11\end{bmatrix}.\]
{Then the Jordan canonical form is given by}
\[A=ZJZ^{-1},\]
{where}
\[Z^{-1}AZ=J=\begin{bmatrix}2&1&0\\0&2&1\\0&0&2\end{bmatrix}\]
{and}
\[Z=\begin{bmatrix}3&0&-1\\0&2&1\\-1&1&1\end{bmatrix},\quad Z^{-1}=\begin{bmatrix}1&-1&2\\-1&2&-3\\2&-3&6\end{bmatrix}.\]
{Using these matrices, $A^{1/2}$ corresponding to $f(x)=x^{1/2}$ is given by}
\[A^{1/2}=Z\begin{bmatrix}\sqrt{2}&\frac{1}{2\sqrt{2}}&-\frac{1}{16\sqrt{2}}\\0&\sqrt{2}&\frac{1}{2\sqrt{2}}\\0&0&\sqrt{2}\end{bmatrix}Z^{-1}.\]
{It is easy to see that $A^{1/2}A^{1/2}=A$.}

{The last example is a matrix exponential function of $A$. The matrix exponential function $e^A$ corresponding to $f(x)=e^x$ is given by}
\[e^A=Z\begin{bmatrix}e^2&e^2&\frac{e^2}{2}\\0&e^2&e^2\\0&0&e^2\end{bmatrix}Z^{-1}.\]
{Notice that computing Jordan canonical form is numerically unstable. These examples, therefore, are not for practical computations.}
\end{Example}

% block: matrix-function-properties-transition page:205 type:ordinary-paragraph
{We enumerate some properties of matrix functions. For details see, e.g., \BookCite{97}.}

% block: matrix-function-properties-theorem pages:205,206 type:theorem
\Needspace{14\baselineskip}
\begin{Theorem}{{Properties of matrix functions}}\label{thm:matrix-function-properties}
{Let $f(x)$ be a function on the spectrum of matrix $A$. Then:}
\begin{enumerate}[label=(\arabic*)]
\item $Af(A)=f(A)A$.
\item $XA=AX\Rightarrow Xf(A)=f(A)X$.
\item $f(A^{\top})=f(A)^{\top}$.
\item $f(XAX^{-1})=Xf(A)X^{-1}$.
\item {If $\lambda_k$ is an eivenvalue of $A$, then $f(\lambda_k)$ is an eigenvalue of $f(A)$.}
\end{enumerate}
{Let $f(x)$ and $g(x)$ be functions on the spectrum of matrix $A$. Then:}
\begin{enumerate}[label=(\arabic*),start=6]
\item {If $(f+g)(x):=f(x)+g(x)$, then $(f+g)(A)=f(A)+g(A)$.}
\item {If $(fg)(x):=f(x)g(x)$, then $(fg)(A)=f(A)g(A)$.}
\end{enumerate}
{Using Definition~\ref{def:function-on-matrix-spectrum}, it follows that}
\begin{enumerate}[label=(\arabic*),start=8]
\item {$f(A)=g(A)$ if and only if $f(\lambda_i)=g(\lambda_i),\ldots,f^{(n_i-1)}(\lambda_i)=g^{(n_i-1)}(\lambda_i)$ $(i=1,2,\ldots,s)$.}
\end{enumerate}
\end{Theorem}

% block: jordan-matrix-function-computational-warning page:206 type:ordinary-paragraph
{From Definition~\ref{def:jordan-matrix-function}, we can define $A^{1/n}$, $e^A$, $\sin A$, $\cos A$, $\log A$ that correspond to elementary functions $x^{1/n}$, $e^x$, $\sin x$, $\cos x$, $\log x$. However, since numerically computing the Jordan canonical form is quite unstable, it is not recommended to compute matrix functions via the Jordan canonical form. Therefore, many researchers have devised numerically stable algorithms for matrix functions.}

% block: matrix-functions-numerical-algorithms-transition page:206 type:ordinary-paragraph
{In what follows, some numerical algorithms are described for computing matrix functions. We will see that the notion of a Krylov subspace is useful for computing some of the matrix functions.}
